\section{The Abel-Jacobi Map}
\label{sec:abel-jacobi}

Fix a base point \(P\in X\). For \(Q\in X\), choose a path from \(P\) to \(Q\)
and define
\[
  I_{P,Q}(\omega)=\int_P^Q\omega.
\]
Changing the path changes \(I_{P,Q}\) by a period, so its class in
\(\Jac(X)\) is well-defined. This is Definition~\ref{def:abel-jacobi}.

\begin{lemma}[Path independence modulo periods]
\label{lem:aj-path-independence}
\uses{def:period-pairing}
The class of \(I_{P,Q}\) in \(\HdualX/\periods_X\) is independent of the path
from \(P\) to \(Q\).
\lean{JacobianChallenge.AbelJacobi.analyticOfCurve}
\leanok
\end{lemma}

\begin{proof}
Two paths from \(P\) to \(Q\) differ by a closed one-cycle. Integrating a
holomorphic one-form over that cycle gives an element of the period lattice.
\end{proof}

\begin{lemma}[Base point maps to zero]
\label{lem:aj-self}
\uses{def:abel-jacobi}
\(\AJ_P(P)=0\).
\lean{Jacobian.ofCurve_self}
\leanok
\end{lemma}

\begin{proof}
Use the constant path from \(P\) to \(P\). The integral of every one-form over
the constant path is zero.
\end{proof}

\begin{lemma}[Holomorphicity of Abel-Jacobi]
\label{lem:aj-holomorphic}
\uses{def:abel-jacobi,prop:complex-torus-package}
The map \(\AJ_P:X\to\Jac(X)\) is holomorphic.
\lean{Jacobian.ofCurve_contMDiff}
\leanok
\end{lemma}

\begin{proof}
Locally, choose a coordinate disk and write a holomorphic one-form as
\(h(z)\,dz\). The coordinate expression for \(\AJ_P\) differs from a constant
by primitives of the holomorphic coefficient functions \(h\). Hence each
coordinate of the lift to \(\HdualX\) is holomorphic, and the quotient map to
\(\Jac(X)\) is holomorphic.
\end{proof}

\subsection{Abel's existence theorem and its consequences}

The named "Abel's theorem" used in
Theorem~\ref{thm:abel-jacobi-injective} is the existence direction:
divisors of degree zero with zero Abel-Jacobi image are principal. We
decompose this and the corollary \(g(X)>0\) ⟹ injectivity into named
sub-obligations so the dep graph reflects the actual classical inputs.

\begin{theorem}[Abel-Jacobi sends divisors to a homomorphism]
\label{thm:aj-divisor-hom}
\uses{def:abel-jacobi}
The Abel-Jacobi map extends \(\Z\)-linearly to a group homomorphism
\(\AJ : \Div^{0}(X) \to \Jac(X)\) on degree-zero divisors. In
particular, \(\AJ_P(Q_1)=\AJ_P(Q_2)\) iff \(\AJ(Q_1-Q_2)=0\).
\depends{Linear extension of \(\AJ_P\) to formal sums; degree-zero
condition is needed to make the result independent of base point.}
\notready
\end{theorem}

\begin{theorem}[Abel existence: zero AJ image implies principal]
\label{thm:abel-existence}
\uses{thm:aj-divisor-hom,def:principal-divisor,input:divisors,input:riemann-bilinear}
For \(D\in\Div^{0}(X)\), if \(\AJ(D)=0\) in \(\Jac(X)\), then there is a
nonzero meromorphic function \(f\in\Mer(X)^{\times}\) with \((f)=D\).
\depends{ABSENT in Mathlib. Classical proof uses the surjectivity of the
exponential of the period map together with a meromorphic-function
construction (theta-function-style). Multi-thousand-line classical
result; see e.g.\ Forster Lectures on Riemann Surfaces, Chapter 3.}
\notready
\end{theorem}

\begin{lemma}[Principal divisor of degree zero with simple two-point support yields a degree-one map]
\label{lem:principal-deg0-simple-support-deg1}
\uses{def:principal-divisor,thm:principal-degree-zero,def:meromorphic-to-cp1}
If \(f\in\Mer(X)^{\times}\) has \((f)=Q_1-Q_2\) with \(Q_1\ne Q_2\), then
the associated map \(\widehat{f}:X\to\CPone\) is nonconstant of
degree~\(1\).
\notready
\end{lemma}

\begin{proofsketch}
\((f)=Q_1-Q_2\) means \(f\) has a simple zero at \(Q_1\), a simple pole
at \(Q_2\), and no other zeros or poles. The pole gives \(\widehat
f^{-1}(\infty)=\{Q_2\}\) with multiplicity 1, so the generic fiber has
size~1.
\end{proofsketch}

\begin{theorem}[Riemann-Hurwitz for degree-1 map: target genus 0]
\label{thm:riemann-hurwitz-deg1}
\uses{def:branched-degree,input:degree-one-isomorphism}
A nonconstant holomorphic map \(f:X\to\CPone\) of degree \(1\) is a
biholomorphism. In particular, \(g(X)=g(\CPone)=0\).
\notready
\end{theorem}

\begin{proof}
Biholomorphism: Input~\ref{input:degree-one-isomorphism} (degree-one
maps are isomorphisms). Genus equality: biholomorphic Riemann surfaces
have isomorphic spaces of holomorphic 1-forms, so the same analytic
genus, and \(g(\CPone)=0\) is a separate well-known computation
(\(H^{0}(\CPone,\Omega^{1})=0\) since the only meromorphic 1-form on
\(\CPone\) up to scaling has poles).
\end{proof}

\begin{theorem}[Point-separation form of Abel's theorem]
\label{thm:abel-point-separation}
\uses{thm:aj-divisor-hom,thm:abel-existence,lem:principal-deg0-simple-support-deg1,thm:riemann-hurwitz-deg1}
If \(g(X)>0\) and \(\AJ_P(Q_1)=\AJ_P(Q_2)\), then \(Q_1=Q_2\).
\notready
\end{theorem}

\begin{proof}
Assume \(Q_1\ne Q_2\). Theorem~\ref{thm:aj-divisor-hom} reduces to
\(\AJ(Q_1-Q_2)=0\). Theorem~\ref{thm:abel-existence} produces a
meromorphic \(f\) with \((f)=Q_1-Q_2\).
Lemma~\ref{lem:principal-deg0-simple-support-deg1} promotes \(f\) to a
degree-one holomorphic \(\widehat f:X\to\CPone\), and
Theorem~\ref{thm:riemann-hurwitz-deg1} forces \(g(X)=0\), contradicting
the hypothesis.
\end{proof}

\begin{classicalinput}[Abel's theorem (umbrella)]
\label{input:abel-theorem}
\uses{thm:abel-existence,thm:abel-point-separation}
\notready
For degree-zero divisors on \(X\), the Abel-Jacobi map has kernel exactly the
principal divisors. In particular, if \(g(X)>0\) and
\(\AJ_P(Q_1)=\AJ_P(Q_2)\), then \(Q_1=Q_2\).
\lean{JacobianChallenge.AbelJacobi.pathIntegralFunctional_separates_points}
\end{classicalinput}

\begin{proof}[Proof of Theorem~\ref{thm:abel-jacobi-injective}]
Assume \(g(X)>0\) and \(\AJ_P(Q_1)=\AJ_P(Q_2)\). By
Theorem~\ref{thm:abel-point-separation}, \(Q_1=Q_2\). Thus \(\AJ_P\) is
injective.
\end{proof}
